% Artifact Appendix template for the NDSS Artifact Evaluation
% version 1.1 (20250525)

% remove the following block when merging the appendix with the camera-ready full paper
%%%
\documentclass[conference]{IEEEtran}
\pagestyle{plain}
\usepackage{url}
\begin{document}
%%%

\appendices
\section{Artifact Appendix}

% The artifact appendix is meant to be a self-contained document presenting a roadmap for setting up and evaluating your artifact. It should provide the following elements:
% \begin{enumerate}
% \item a list of the hardware, software, and configuration \textbf{requirements} for running the artifact;
% \item a clear description of how, and in what respects, the artifact \textbf{supports} the research presented in the paper;
% \item a guide for how others can \textbf{execute and validate} the artifact for its functional and usability aspects;
% \item the \textbf{major claims} made by your paper and a clear description of how to obtain data for each claim through your supplied artifact.
% \end{enumerate}

% Points 3-4 are not required if only the \textit{Available} badge is requested. Linking the paper claims to the artifact is a necessary step that allows artifact evaluators to functionally test and reproduce your results. Towards that end, explicitly list down items (e.g., results, plots, tables) from the paper and cross-reference them with the experiments to be reproduced.
 
% Unless otherwise stated, filling every section below with the requested contents is mandatory for participating in the artifact evaluation process. Section titles cannot be changed.

\subsection{Description \& Requirements}

% This section should list all the information necessary to recreate the experimental setup you used to run your artifact.

% Provide also a link to an archival repository where all the artifact's main components (software, datasets, documentation, etc.) can be accessed and, where this applies, the minimal hardware and software requirements to run your artifact.

% It is also very good practice to list and describe in this section benchmarks where those are part of your artifacts or simply have been used to produce results with it.

The artifact contains all the code and binaries to reproduce the paper's claims. 
The mitigation is stored in \texttt{optee\_shm\_patch}. \texttt{ae} contains the artifact evaluation wrapper scripts.
The remaining folders all pertain to ScHMuzz. 

\subsubsection{How to access}
The artifact can be downloaded via hotcrp.
% Describe here how to access your artifact. During the artifact evaluation, in case of a private repository, you should provide instructions on how to access it. For the camera-ready version of the appendix, you must provide a DOI link to the AEC-approved artifact version that you must upload by then to permanent storage.

\subsubsection{Hardware dependencies}
x86-64 CPU with 8 cores and 16 GB of RAM running a recent Linux.
% Simply write ``None." where this does not apply to your artifact.

\subsubsection{Software dependencies}
Docker is required to run the artifact.

\subsubsection{Benchmarks}
None
% Describe here any data (e.g., datasets, models, workloads, etc.) required by the experiments with this artifact reported in your paper. Simply write ``None." where this does not apply to your artifact.


\subsection{Artifact Installation \& Configuration}

Download and unpack the artifact archive. 

Then run the following command to setup: 

\texttt{./ae/ae.sh setup}

% This section should include all the high-level installation and configuration steps required to prepare the environment to be used for the evaluation of your artifact.

To run experiement E4 please send us your SSH public key over hotcrp. 
We will add it to the authorized keys for our VPS. The VPS 
will be used to access our testing devices for E4 over adb. 

\subsection{Experiment Workflow}

All experiments except E4 are run from the \texttt{ae/ae.sh} scripts. 
The script can be used to run individual experiments or all relevant sequentially. 
% This section should provide a high-level view of the experimental workflow and how it is implemented, invoked, and (if needed) customized. The section is optional if the experiment workflow can be easily embedded in the Evaluation section.

For E4, first connect with SSH to our VPS: 

\small
\texttt{ssh ae@65.108.89.50}
\normalsize 

Then follow the instructions for the experiment.
% \textit{Detailed information for submission, reviewing, and badging process followed for the evaluation of NDSS 2027 artifacts can be found at:} \url{https://secartifacts.github.io/ndss2027/}.

\subsection{Major Claims}
% Enumerate here the major claims (Cx) made in your paper. For this purpose, we ask you to use the this and the following section and cross-reference the items therein, as explained next.

% Follows an example:

\begin{itemize}
    \item (C1): ScHMuzz automatically detects crashes only triggerable via double fetches in TAs (Table 1, Figure 4, Figure 5). This is proven by experiment E1
    \item (C2): We found six double fetch vulnerabilities (Table 2). This is proven by experiment E2.
    \item (C3): Rust TAs are also affected by overlapping fetches (Table 4). This is proven by experiment E3.
    \item (C4): Teegris, Mitee, QSEE, and Kinibi support zero-copy shared memory (Table 3). This is proven by experiment E4.
\end{itemize}

\subsection{Evaluation}

% This section include all the operational steps and experiments which must be performed to evaluate if your artifact is functional and validate the proposed/presented experiments. For this purpose, we ask you to use the this and the preceding section and cross-reference the items therein.

% If the \textit{Reproduced} badge is requested, the section should reference your paper's key results: describe the expected numbers and, where applicable, the maximum variation expected (particularly important for performance numbers). If an experiment is a scaled-down version of the one run for the paper's evaluation, provide a clear indication of it with a justification of the scaling-down being meaningful.

% We also highly encourage you to provide your estimates of human- and compute-time for each of the listed experiments.

% Follows an exemplary structure for one experiment (Ey):

\subsubsection{Experiment (E1)}
[Automatic Double Fetch Detection] [45 compute-hour + 10 human-minutes]: Runs the fuzzing campaign on the 30 harnessesd TAs including all three of ScHmuzz's stages.

% \textit{[How to]} 
% Describe thoroughly the steps to perform the experiment and to collect and organize the results as expected from your paper. We encourage you to use the following structure with three blocks for the description of an experiment (and to provide precise indications about the expected outcome for each of the steps from the blocks):

\textit{[Preparation]}
\small{\texttt{./ae/ae.sh setup}}
\normalsize

\textit{[Execution]}
\footnotesize
\texttt{./ae/ae.sh e1\_automatic\_df\_detection}
\normalsize

\textit{[Results]}
The script prints the resulting data in the same format as Table 1 and generates the coverage graph figures (Figure 4,5).
Compare the results from the table to Table 1 in the paper, the numbers should be similar. Compare the coverage graphs with 
Figure 4 and 5 from the paper, again they should be similar. By default this eval reduces fuzzing times to allow it to run on a commodity machine, the numbers will not be the same as in the paper.

\subsubsection{Experiment (E2)}
[TOCTTOU Vulnerabilities] [\textless1 compute-hour + 5 human-minutes]: Runs the proof of concepts for the six discovered vulnerabilities against the emulator unilt the crashes reproduce.

% \textit{[How to]} 
% Describe thoroughly the steps to perform the experiment and to collect and organize the results as expected from your paper. We encourage you to use the following structure with three blocks for the description of an experiment (and to provide precise indications about the expected outcome for each of the steps from the blocks):

\textit{[Preparation]}
None

\textit{[Execution]}
\small
\texttt{./ae/ae.sh e2\_vulns}
\normalsize

\textit{[Results]}
The script runs until it triggers the crashes for each vulnerability. Output should be that all the vulnerabilities reproduce.

\subsubsection{Experiment (E3)}
[Rust TAs] [3 compute-hours + 5 human-minutes]: Runs ScHMuzz on the five Rust TAs.

% \textit{[How to]} 
% Describe thoroughly the steps to perform the experiment and to collect and organize the results as expected from your paper. We encourage you to use the following structure with three blocks for the description of an experiment (and to provide precise indications about the expected outcome for each of the steps from the blocks):

\textit{[Preparation]}
None

\textit{[Execution]}
\small
\texttt{./ae/ae.sh e3\_rust}
\normalsize

\textit{[Results]}
After finishing the script prints the data in the format of Table 4. Compare the results, for 3 out of the 5 TAs ScHMuzz should detect overlapping fetches.

\subsubsection{Experiment (E4)}
[On-device Shared Memory] [\textless1 compute-hour + 5 human-minutes]: runs pocs on-device that trigger double fetches to confirm zero-copy shared memory.

% \textit{[How to]} 
% Describe thoroughly the steps to perform the experiment and to collect and organize the results as expected from your paper. We encourage you to use the following structure with three blocks for the description of an experiment (and to provide precise indications about the expected outcome for each of the steps from the blocks):

\textit{[Preparation]}
Obtain access to the VPS and login via SSH.

\textit{[Execution]}
\small
\texttt{cd oversharing\_ae}

\texttt{./ae/ae\_ondevice.sh}
\normalsize

\textit{[Results]}
The script prints the results in the same format as Table 3.

\subsection{Customization}
The \texttt{ae.sh} script can be customized using a number of environment varaiables. 

\texttt{AE\_SCALE=paper} uses the same times as the paper's evaluation.

\texttt{AE\_JOBS} scales the number of parallel emulator containers.

\texttt{AE\_EXPLORE\_TIME} scales the time to fuzz each TA in the exploration phase.

\texttt{AE\_FAF\_TIME} sclaes the time spent on fetch anchored fuzzing per snapshot.

% Provide here notes on how to customize your experiments, when applicable. The section is optional.

\subsection{Notes}
% This section is meant to allow authors to include any further important notes that may not fall within any of the previous categories. We kindly encourage you to remove this section where this sort of content may not be strictly needed (rather than filling it with unnecessary or redundant information).
The reproduction for C5 includes only a subset of the testing devices used for the experiment in the paper. 
This is because some of these devices are currently being used for other projects. Additionally, both our mitee testing devices were unrooted during 
an update and can not be easily rerooted.

% remove the following block when merging the appendix with the camera-ready full paper
%%%
\end{document}
%%%
